\section{Limitations}

Although GLM-OCR demonstrates competitive performance across diverse benchmarks and practical scenarios, several limitations remain.

\subsection{Two-Stage Architectural Constraints}

The current two-stage pipeline, consisting of layout analysis followed by region-level recognition, may introduce error propagation. In cases of inaccurate layout detection, downstream recognition performance may degrade. Additionally, complex layouts involving cross-page dependencies or irregular multi-column structures may lead to imperfect reading order reconstruction.

\subsection{Data Coverage Limitations}

Model performance is influenced by the distribution and diversity of training data. Degradation may occur in scenarios involving:

\begin{itemize}
    \item Extremely low-resolution or heavily distorted documents,
    \item Highly complex mathematical expressions,
    \item Dense or irregular tabular structures,
    \item Underrepresented languages in the training corpus.
\end{itemize}

\subsection{Structured Output Variability}

As a generative model, GLM-OCR may exhibit minor stochastic variation in formatting behaviors, particularly in line breaks and whitespace handling. Although reinforcement learning and structural supervision mitigate this effect, strict formatting guarantees cannot be fully ensured.

\subsection{Key Information Extraction}

While the model supports prompt-based KIE, extraction accuracy depends on prompt specification and schema clarity. In complex forms with implicit or ambiguous field boundaries, incomplete or redundant outputs may occur.

These limitations represent areas for continued research and system refinement.